% answers_short.tex —— 2004 年·普通高等学校春季招生考试·安徽卷理科 · 简短答案速查
% 由 tools/make_short_answers.py 自 answers.tex 的 \gksolhead/\gkanswer 提取，
% 并与 2004/2004_数学真题_春季安徽理_答案要点.md 逐题比对后生成，勿手改（改源 MD 后重跑）。

\begin{gknotice}
\textbf{使用说明}：本节只给最终答案便于快速核对，完整推导见解析版。答案为编写者独立推导，\textbf{非官方标准答案}。
\end{gknotice}

\gkpart{选择题}
\begin{center}\begin{tabular}{ccccccccccccc}
\toprule
题号 & 1 & 2 & 3 & 4 & 5 & 6 & 7 & 8 & 9 & 10 & 11 & 12 \\
\midrule
答案 & D & A & C & C & A & C & D & D & D & B & C & B \\
\bottomrule
\end{tabular}\end{center}

\gkpart{填空题}
\begin{itemize}[leftmargin=2.6em]
  \item[\textbf{13.}] $x=-\dfrac32$
  \item[\textbf{14.}] $\dfrac7{10}$
  \item[\textbf{15.}] $\dfrac14$
  \item[\textbf{16.}] $n=3$
\end{itemize}

\gkpart{解答题}
\begin{itemize}[leftmargin=2.6em]
  \item[\textbf{17.}] 若 $a>1$，解集为 $(0,a^{-\sqrt3})\cup(1,a^{\sqrt3})$；若 $0<a<1$，解集为 $(a^{\sqrt3},1)\cup(a^{-\sqrt3},+\infty)$
  \item[\textbf{18.}] $b_n=2n-1$
  \item[\textbf{19.}] (1) $k=1$ 时为圆，$k\ne1$ 时为椭圆；(2) $\max(|y|,k|x|)=\dfrac{c\sqrt{1+k^2}}2$，轨迹为矩形周界
  \item[\textbf{20.}] (1) $\cos\angle(AC,BC_1)=\dfrac{\sqrt{10}}5$；(2) $A_1B\perp$ 面 $AB_1C$
  \item[\textbf{21.}] $P(\xi=2)=\dfrac{28}{45}$，$P(\xi=3)=\dfrac{14}{45}$，$P(\xi=4)=\frac{1}{15}$；$E\xi=\dfrac{22}{9}$
  \item[\textbf{22.}] (1) $M=(-1,-\dfrac{19}{7})$；(2) 对任意 $a$，顶点 $M=(-2,-\dfrac{26}{7})$ 处的竖直法线 $x=-2$ 均通过 $P$。若 $a>-\dfrac{45}{14}$，另有两点 $M=(-2\pm\sqrt{a+\dfrac{45}{14}},a-\dfrac12)$，法线为 $y-a=\mp\dfrac{x+2}{2\sqrt{a+\dfrac{45}{14}}}$；若 $a\le-\dfrac{45}{14}$，没有其他点
\end{itemize}
